\originalchapter{含时微扰增补：连续谱、光电离与辐射}
本章接第8章的动力学表象和第15章的含时跃迁, 根据Barton Zwiebach的MIT OCW 8.06 Spring 2018 Chapter 4, 原件印刷页73--110重编. 参考或改编署名不使整章失去OCW默认许可. 下述说明和推导为中文独立整理; 原件有两处符号笔误在对应计算中按物理条件纠正. 本章只采用非相对论量子力学、弱外场、电偶极和注明的连续末态近似.
\originalsection{瞬时本征态、脉冲和微扰阶数}
含时Hamilton算符在每个固定时刻可以有瞬时本征态; 一般没有一个固定空间波函数乘$\e^{-\ii Et/\hbar}$的全程定态. 这两句话并不相同. 演化仍可用$H_0$的完整本征基展开, 时间变化交给系数. 对初态$\ket i$, 一阶跃迁幅为第15章的时间积分. 若$W(t)$厄米且从同一$t_0$到$t_1$施加, $\omega_{if}=-\omega_{fi}$和$W_{if}=W_{fi}^*$使$c_{i\leftarrow f}^{(1)}=-[c_{f\leftarrow i}^{(1)}]^*$, 所以两个方向在这一最低非零概率阶相等. 这不是任意强驱动的精确详细平衡定律.

逐阶展开$\psi_I=\sum_{k\ge0}\lambda^k\psi_I^{(k)}$, 令初态与$\lambda$无关, 则$\psi_I^{(k>0)}(0)=0$, $\ii\hbar\partial_t\psi_I^{(k)}=W_I(t)\psi_I^{(k-1)}$. 反复积分自动给时间有序区域, 与第8章的二阶公式相接. 对已有叠加初态$\sum c_n(0)\ket n$, 一阶校正为
\[
c_m^{(1)}(t)=-\frac{\ii}{\hbar}\sum_n c_n(0)
\int_0^t W_{mn}(s)\e^{\ii\omega_{mn}s}\dd s.
\]
若$c_m(0)\ne0$, 一阶概率是$|c_m(0)|^2+2\operatorname{Re}[c_m(0)^*c_m^{(1)}]$. 单独保留$|c_m^{(1)}|^2$而丢$c_m^{(2)}$与$c_m(0)$的干涉项, 不能宣称概率精确到二阶. 从一个本征态转到正交态时零阶振幅为0, 才能用一阶幅平方给最低二阶概率.

一个能看清强弱边界的例子是瞬时脉冲$H=H_0+K\delta(t)$, $K=\left(\begin{smallmatrix}0&\alpha\\\alpha^*&0\end{smallmatrix}\right)$, $\alpha$量纲为作用量, 原件注明例子based on Griffiths 9.3. 把$\delta$改为宽$\epsilon$、高$1/\epsilon$的脉冲, 脉冲演化为$\exp[-\ii(H_0\epsilon+K)/\hbar]$. 有限维取$\epsilon\to0$给$U_{\rm kick}=\e^{-\ii K/\hbar}$; 不能直接将不连续态在$\delta$处取一个任意端点值. 因$K^2=|\alpha|^2I$,
\[
U_{\rm kick}=I\cos\frac{|\alpha|}{\hbar}
-\ii\frac{K}{|\alpha|}\sin\frac{|\alpha|}{\hbar}.
\]
从$\ket a$出发, $\ket{\psi(0+)}=\cos(|\alpha|/\hbar)\ket a-\ii(\alpha^*/|\alpha|)\sin(|\alpha|/\hbar)\ket b$, $P_b=\sin^2(|\alpha|/\hbar)$. $\alpha=0$取连续极限. 弱脉冲才可展开成$P_b\simeq|\alpha|^2/\hbar^2$; 时间很短不等于总作用很小.
\originalsection{磁共振的精确校验}
设$H_0=\omega_0S_z$, 先在$t=0$加定常$\Omega S_x$. 总场方向为$(\Omega,0,\omega_0)$, 模长$\nu=\sqrt{\omega_0^2+\Omega^2}$. 利用$(\boldsymbol n\cdot\boldsymbol\sigma)^2=I$得到从$z$向上自旋到向下的精确概率$P_-(t)=(\Omega^2/\nu^2)\sin^2(\nu t/2)$. 小$\Omega/\omega_0$时最低阶为$(\Omega/\omega_0)^2\sin^2(\omega_0t/2)$; 对非常长时间, 频率的小修正也可累计成为有限相位误差, 不能只凭小振幅声称所有时间的相对误差都小.

若加旋转场$W(t)=\Omega(S_x\cos\omega_0t+S_y\sin\omega_0t)$, 对易关系给$\e^{\ii\omega_0tS_z/\hbar}S_x\e^{-\ii\omega_0tS_z/\hbar}=S_x\cos\omega_0t-S_y\sin\omega_0t$. 因而相互作用表象中$W_I=\Omega S_x$, 正好定常. 精确演化为$U=\e^{-\ii\omega_0tS_z/\hbar}\e^{-\ii\Omega tS_x/\hbar}$, 从$z$向上出发$P_-=\sin^2(\Omega t/2)$, 自旋可以完全翻转. 此处旋转方向和共振频率均由算符共轭实际确定, 不靠忽略一个快振荡项得到精确性.
\originalsection{连续末态计数与黄金规则}
将空间先放在边长$L$的周期盒内, 末态$u_{\boldsymbol k}=L^{-3/2}\e^{\ii\boldsymbol k\cdot\boldsymbol r}$, $k_j=2\pi n_j/L$. 一个$\dd^3k$小格中态数是$L^3\dd^3k/(2\pi)^3$. 用$E=\hbar^2k^2/(2m)$, $k\dd k=m\dd E/\hbar^2$, 得每单位能量、单位立体角的态密度
\[
\rho_{E,\Omega}=\frac{L^3mk}{8\pi^3\hbar^2},\qquad
\sum_fF_f\longrightarrow\frac{L^3}{(2\pi)^3}\int F(\boldsymbol k)\dd^3k.
\]
平面波只在自由或末态相互作用可忽略时合适. 局域初态到一个盒末态的矩阵元含$L^{-3/2}$, 平方与态密度的$L^3$抵消, 预测不应依赖人为盒长.

定常弱扰动从0开到$t$时,$P_{fi}=4|W_{fi}|^2\sin^2[(E_f-E_i)t/(2\hbar)]/(E_f-E_i)^2$. 对连续末态求和并在窄窗近似固定矩阵元、密度. 所用核面积可以直接证明: $\sin u/u=\tfrac12\int_{-1}^1\e^{\ii su}\dd s$, 按Parseval等式（傅里叶变换规范$\hat f(u)=\int f(s)\e^{\ii us}\dd s$）, $\int_{\R}(\sin u/u)^2\dd u=\pi$. 变量代换$u=(E_f-E_i)t/(2\hbar)$于是给总概率$2\pi t|W_{fi}|^2\rho_E(E_i)/\hbar$. 一方面$t$足够大使共振窗宽$\sim\hbar/t$内谱与耦合变化小, 另一方面$\Gamma t\ll1$避免初态耗尽, 在两者共同窗口才有$\Gamma=2\pi|W_{fi}|^2\rho_E(E_i)/\hbar$.

对$W(t)=2H'\cos\omega t$, 两个Fourier分量的幅度均为$H'$, 不能再把峰值$2H'$平方代进黄金规则. 有限时幅为
\[
c_f^{(1)}=-\frac{H'_{fi}}{\hbar}\left[
\frac{\e^{\ii(\omega_{fi}+\omega)t}-1}{\omega_{fi}+\omega}
+\frac{\e^{\ii(\omega_{fi}-\omega)t}-1}{\omega_{fi}-\omega}\right].
\]
分母零时取连续极限. 吸收共振$E_f=E_i+\hbar\omega$, 受激发射$E_f=E_i-\hbar\omega$, 两者率按各自末态密度计算. 外加经典周期场得到的是受激发射; 自发发射还需真空辐射模式, 或下文的热平衡论证. 原件在此局部把受激发射写成自发发射, 中文按这一区别修正.
\originalsection{氦的自电离}
忽略电子间排斥的氦模型必须是$H_0=\sum_{j=1}^2[p_j^2/(2m)-Ze^2/r_j]$, $Z=2$. 原件一行漏写$Z$, 其后能量含$Z^2$, 本处保持两者一致. 能量$E_{n_1n_2}=-Z^2E_{\rm Ry}(n_1^{-2}+n_2^{-2})$. 一个电子留在$n$束缚态、另一个电子自由的连续通道从$-Z^2E_{\rm Ry}/n^2$开始; 总能负不表示两个电子一定都束缚.

例如两个2S电子的未扰能$-2E_{\rm Ry}=-27.2$ eV, 已高于1S加自由电子通道阈值$-4E_{\rm Ry}=-54.4$ eV. 扰动$W=e^2/|\boldsymbol r_1-\boldsymbol r_2|$可把一个电子降入1S, 能量释放给另一个电子, 使其带$2E_{\rm Ry}=27.2$ eV动能离开. 没有外来光子仍可发生电离, 能量在电子间交换. 计算率应将正确交换对称性的初末态代入$W$矩阵元并积分连续末态; 本例只说明能量通道及黄金规则对象, 没有宣称算出真实氦共振寿命.
\originalsection{氢光电离的矩阵元与角分布}
采用原件的Gaussian电磁单位, $e>0$为电荷大小, 氢初态$\psi_{1s}=\e^{-r/a_0}/\sqrt{\pi a_0^3}$, $E_i=-E_{\rm Ry}$. 光场$\boldsymbol E=2E_0\cos\omega t\,\hat z$, 电偶极扰动$2eE_0z\cos\omega t$, 故$H'=eE_0z$. 光子波长远大于原子尺度时可忽略场的空间变化, 需$k_\gamma a_0\ll1$; 用末态平面波还需电子动能远大于库仑尺度$E_{\rm Ry}$. 可同时取$E_{\rm Ry}\ll\hbar\omega\ll\hbar c/a_0=2E_{\rm Ry}/\alpha$, 并保持$\hbar\omega\ll mc^2$. 仅$\lambda\gg a_0$较松的条件并不精确控制$k_\gamma a_0$的误差; 这里给出明确长波小量. 阈值附近需库仑连续波函数, 不能把下述平面波公式当精确阈值定律.

令$F(\boldsymbol k)=\int\e^{-r/a_0-\ii\boldsymbol k\cdot\boldsymbol r}\dd^3r$. 取极轴沿$\boldsymbol k$, 角积分$4\pi\sin(kr)/(kr)$, 再用$\int_0^\infty r\e^{-r/a_0}\sin kr\dd r=2(k/a_0)/(k^2+a_0^{-2})^2$, 得$F=8\pi a_0^3/(1+k^2a_0^2)^2$. 对$k_z$求导将$z$带入积分, 因$\partial_{k_z}F=-\ii\int z\e^{-r/a_0-\ii\boldsymbol k\cdot\boldsymbol r}\dd^3r$,
\[
H'_{\boldsymbol k,1s}=\frac{\ii eE_0}{\sqrt{\pi a_0^3L^3}}\partial_{k_z}F
=-\frac{32\ii\sqrt\pi\,eE_0a_0^5k\cos\theta}
{\sqrt{a_0^3L^3}(1+k^2a_0^2)^3}.
\]
$\theta$是电子射出方向与电场偏振夹角, 正负相位不影响率. 用前节态密度及$\hbar^2k^2/(2m)=\hbar\omega-E_{\rm Ry}$, 设$q=ka_0=\sqrt{\hbar\omega/E_{\rm Ry}-1}$, 得
\begin{align*}
\frac{\mathrm d\Gamma}{\mathrm d\Omega}
&=\frac{256}{\pi}\frac{ma_0^2}{\hbar^2}\frac{(eE_0a_0)^2}{\hbar}
\frac{q^3\cos^2\theta}{(1+q^2)^6},\\
\Gamma&=\frac{512}{3}\frac{(eE_0a_0)^2}{\hbar E_{\rm Ry}}
\frac{q^3}{(1+q^2)^6}.
\end{align*}
角积分用$\int\cos^2\theta\dd\Omega=4\pi/3$, 因而沿偏振方向最强, 正交方向为0. 高能$q\gg1$下总率随$q^{-9}$下降. 若需光电离截面, 再除光子流: Gaussian单位的平均强度$I=cE_0^2/(2\pi)$, 光子流$I/(\hbar\omega)$, 故$\sigma=\Gamma\hbar\omega/I$. 场幅度已在分子分母抵消, 确认截面在弱光线性响应下与强度无关.
\originalsection{Einstein系数和偶极寿命}
设两个不简并能级$a,b$, $\Delta=E_b-E_a=\hbar\omega_0>0$. 热平衡占据比是$N_b/N_a=\e^{-\beta\Delta}$, $\beta=1/(k_BT)$; 原件式4.5.2把占据比左右写反, 后续计算使用的实际是这一正确比值. 设$U(\omega)$为每单位角频率、单位体积的辐射能量密度, 吸收率$B_{ab}U(\omega_0)N_a$, 受激发射率$B_{ba}U(\omega_0)N_b$, 自发发射率$AN_b$. 平衡条件给
\[
U(\omega_0)=\frac{A}{B_{ab}\e^{\beta\Delta}-B_{ba}}.
\]
与Planck密度$U(\omega)=\hbar\omega^3/[\pi^2c^3(\e^{\beta\hbar\omega}-1)]$在所有温度的函数形式比较, 得$B_{ab}=B_{ba}=B$, $A/B=\hbar\omega_0^3/\pi^2c^3$. 简并能级合并统计时改为$g_aB_{ab}=g_bB_{ba}$, 不应直接搬用单态等式.

令偶极矩$\boldsymbol d=q_{\rm charge}\boldsymbol r$, $W=-\boldsymbol d\cdot\boldsymbol E$. 对单偏振单位向量$\boldsymbol n$, 用谐波跃迁核在非相干宽频光谱上积分, $E_0^2=2\pi\langle u\rangle$给概率$P(t)=8\pi|\boldsymbol d_{ab}\cdot\boldsymbol n|^2\hbar^{-2}\int U(\omega)\sin^2[(\omega_0-\omega)t/2]/(\omega_0-\omega)^2\dd\omega$. 窄共振窗把$U$取为$U(\omega_0)$, 令$u=(\omega-\omega_0)t/2$, 积分核面积为$\pi t/2$, 除$t$给单原子率$4\pi^2|\boldsymbol d_{ab}\cdot\boldsymbol n|^2U(\omega_0)/\hbar^2$. 非相干表示不同频率与偏振的随机相位使交叉项平均消失, 因而相加的是概率. 各向同性光场满足$\langle n_in_j\rangle=\delta_{ij}/3$, 这是旋转对称性及$\sum_i n_i^2=1$的结果. 因而
\[
B=\frac{4\pi^2}{3\hbar^2}|\boldsymbol d_{ab}|^2,\qquad
A=\frac{4\omega_0^3}{3\hbar c^3}|\boldsymbol d_{ab}|^2.
\]
这里$|\boldsymbol d|^2=\sum_jd_j^*d_j$, 矩阵元向量可以是复数. 换成SI单位,$A=\omega_0^3|\boldsymbol d_{ab}|^2/(3\pi\epsilon_0\hbar c^3)$. 不把两种单位的$4\pi\epsilon_0$混用. 单通道寿命$\tau=1/A$, 多通道先加率$\tau^{-1}=\sum_aA_{b\to a}$. 在固定率的独立衰减模型下$\dot N_b=-AN_b$给$N_b(t)=N_b(0)\e^{-At}$; 初始极短时间的精确量子衰减仍可为二次, 这一指数是速率近似.
\originalsection{电偶极选择定则}
径向波函数与球谐函数分离时,$z=r\cos\theta$,$x\pm\ii y=r\sin\theta\e^{\pm\ii\phi}$. 对$\phi$积分, 指数只有零总幂才不消失, 所以$z$要求$\Delta m=0$, $x\pm\ii y$要求$\Delta m=\pm1$. 更强的$\ell$条件来自球谐递推: $\cos\theta Y_{\ell m}=c_+Y_{\ell+1,m}+c_-Y_{\ell-1,m}$, 系数$c_+=\sqrt{[(\ell+1)^2-m^2]/[(2\ell+1)(2\ell+3)]}$, $c_-=\sqrt{[\ell^2-m^2]/[(2\ell-1)(2\ell+1)]}$. 该式由关联Legendre多项式三项递推及归一化得到; $\sin\theta\e^{\pm\ii\phi}Y_{\ell m}$同样只含$\ell\pm1,m\pm1$, 具体$[L_\pm,\cos\theta]=\mp\hbar\sin\theta\e^{\pm\ii\phi}$, 用$L_\pm Y_{\ell m}=\hbar\sqrt{(\ell\mp m)(\ell\pm m+1)}Y_{\ell,m\pm1}$作用递推式两侧, 再扣$\cos\theta L_\pm Y_{\ell m}$, 所剩仍只在$\ell\pm1,m\pm1$空间, 得横分量结论. 球谐正交性于是迫使$\Delta\ell=\pm1$. 宇称也验证它: 电偶极为奇宇称算符, 同宇称态的角积分为0.

这些条件只是允许非零的必要条件, 径向积分仍可能消失. 自旋无关的电偶极算符保持自旋标签. 例如2S到1S的单光子电偶极矩阵元为0, 不表示所有更高阶辐射通道都禁绝; 要区分本近似下的选择定则与真实原子的全部寿命机制.
